%---------------------------------------------------------------------
% Brief edition: the sixteen selected questions.
%
% Solutions are compressed relative to the 27-question edition: fewer
% display lines, one \parakh only where a sign error is genuinely likely,
% and only five \ghalti boxes -- the rest are collected on the last page.
%---------------------------------------------------------------------

\section*{یونٹ \textenglish{1}: احتمال}
\markboth{یونٹ 1: احتمال}{}

\begin{nukta}[چار اصول]
\[
\prob{A'}=1-\prob{A},\qquad
\prob{A\cup B}=\prob{A}+\prob{B}-\prob{A\cap B}
\]
\[
\prob{A\cap B}=\prob{A}\prob{B}\ \ (\text{آزاد}),\qquad
\prob{A\mid B}=\frac{\prob{A\cap B}}{\prob{B}}
\]
باہم خارج واقعات میں \(\prob{A\cap B}=0\)، تب جمع کا اصول سادہ ہو جاتا ہے۔
\end{nukta}

\begin{sawal}{سوال \textenglish{1}}
ایک مرتبان میں \textenglish{8} سبز نقطہ دار، \textenglish{10} سبز دھاری دار،
\textenglish{12} نیلی نقطہ دار اور \textenglish{10} نیلی دھاری دار گیندیں ہیں۔
ایک گیند بے ترتیب نکالی جائے تو احتمال کیا ہے کہ وہ \textbf{(الف)} سبز یا دھاری
دار ہو، \textbf{(ب)} نقطہ دار ہو، \textbf{(ج)} نیلی یا نقطہ دار ہو؟
\end{sawal}

\begin{hal}
پہلے دو رخی جدول بنائیے --- اس سوال کے سارے نمبر اسی قدم پر ملتے یا کھوتے ہیں۔

\begin{center}
\begin{tabular}{@{}p{2.4cm} p{2.2cm} p{2.2cm} p{1.8cm}@{}}
\toprule
 & \textbf{نقطہ دار} & \textbf{دھاری دار} & \textbf{کل} \\
\midrule
سبز  & \textenglish{8}  & \textenglish{10} & \textbf{\textenglish{18}} \\
نیلی & \textenglish{12} & \textenglish{10} & \textbf{\textenglish{22}} \\
\midrule
\textbf{کل} & \textbf{\textenglish{20}} & \textbf{\textenglish{20}} & \textbf{\textenglish{40}} \\
\bottomrule
\end{tabular}
\end{center}

\textbf{(الف)} سبز اور دھاری دار میں اشتراک ہے (\textenglish{10} گیندیں دونوں
میں)، اس لیے عام جمع کا اصول:
\[
\prob{G\cup S}=\tfrac{18}{40}+\tfrac{20}{40}-\tfrac{10}{40}=\tfrac{28}{40}=0.70
\]
\textbf{(ب)} \(\prob{D}=\tfrac{20}{40}=0.50\)۔ \quad
\textbf{(ج)} اشتراک \textenglish{12} نیلی نقطہ دار:
\(\prob{B\cup D}=\tfrac{22}{40}+\tfrac{20}{40}-\tfrac{12}{40}=0.75\)

\parakh{براہِ راست گنتی: \(8+10+10=28\) اور \(12+10+8=30\)۔ مطابقت ہے۔}

\jawab{(الف) \(0.70\) \quad (ب) \(0.50\) \quad (ج) \(0.75\)}
\end{hal}

\begin{ghalti}
سادہ جمع کا اصول \(\prob{A}+\prob{B}\) صرف \emph{باہم خارج} واقعات پر چلتا ہے۔
اشتراک والے واقعات پر لگانے سے اشتراک دو بار گنا جاتا ہے --- اس یونٹ کی سب سے
عام غلطی یہی ہے۔
\end{ghalti}

\begin{sawal}{سوال \textenglish{2}}
اگر \(\prob{E}=0.2\)، \(\prob{F}=0.5\) اور \(\prob{E\cup F}=0.6\) ہو تو نکالیے
\textbf{(الف)} \(\prob{E\cap F}\)، \textbf{(ب)} \(\prob{E\mid F}\)،
\textbf{(ج)} \(\prob{F\mid E}\)۔
\end{sawal}

\begin{hal}
\textbf{(الف)} جمع کے اصول کو الٹ کر اشتراک کو مطلوب بنائیے:
\[
\prob{E\cap F}=\prob{E}+\prob{F}-\prob{E\cup F}=0.2+0.5-0.6=0.1
\]
\textbf{(ب)} \(\prob{E\mid F}=\tfrac{0.1}{0.5}=0.2\) \qquad
\textbf{(ج)} \(\prob{F\mid E}=\tfrac{0.1}{0.2}=0.5\)

\medskip
\textbf{اضافی نمبر یہاں ملتے ہیں:} \(\prob{E\mid F}=0.2=\prob{E}\) اور
\(\prob{E}\prob{F}=(0.2)(0.5)=0.1=\prob{E\cap F}\)۔ لہٰذا \(E\) اور \(F\)
\textbf{آزاد} واقعات ہیں۔

\jawab{(الف) \(0.1\) \quad (ب) \(0.2\) \quad (ج) \(0.5\)؛ \(E,F\) آزاد ہیں۔}
\end{hal}

%---------------------------------------------------------------------
\section*{یونٹ \textenglish{2}: بے ترتیب متغیر}
\markboth{یونٹ 2: بے ترتیب متغیر}{}

\begin{sawal}{سوال \textenglish{3}}
\textbf{(الف)} ایک کھرا سکہ تین بار اچھالا جاتا ہے؛ \(X\) چِتوں کی تعداد ہے۔
تقسیم بنائیے اور \(P(X\ge 2)\) نکالیے۔
\textbf{(ب)} اگر \(X=0,2,4,6,8\) کے احتمالات \(K,\ 2K,\ 4K,\ 2K,\ K\) ہوں تو
\(K\) نکالیے۔
\textbf{(ج)} ایک تھیلے میں \textenglish{4} سرخ اور \textenglish{6} کالی گیندیں
ہیں۔ \emph{بغیر واپسی} کے \textenglish{4} گیندیں نکالی جاتی ہیں؛ \(X\) سرخ
گیندوں کی تعداد ہے۔ \(P(X=2)\) نکالیے۔
\end{sawal}

\begin{hal}
\textbf{(الف)} کل \(2^{3}=8\) یکساں ممکن نتائج ہیں:

\begin{center}
% The last column has to hold "P(X=x)" unbroken; at 1.8cm it wrapped
% mid-expression.
\begin{tabular}{@{}p{1.3cm} p{4.6cm} p{1.5cm} p{2.4cm}@{}}
\toprule
\(x\) & نتائج & تعداد & \(P(X=x)\) \\
\midrule
\textenglish{0} & \textenglish{TTT}           & \textenglish{1} & \(1/8\) \\
\textenglish{1} & \textenglish{HTT, THT, TTH} & \textenglish{3} & \(3/8\) \\
\textenglish{2} & \textenglish{HHT, HTH, THH} & \textenglish{3} & \(3/8\) \\
\textenglish{3} & \textenglish{HHH}           & \textenglish{1} & \(1/8\) \\
\bottomrule
\end{tabular}
\end{center}

''دو یا زیادہ'' کا مطلب \(X=2\) یا \(X=3\)، اور یہ باہم خارج ہیں:
\[
P(X\ge 2)=\tfrac38+\tfrac18=\tfrac12=0.5
\]

\textbf{(ب)} احتمالات کا مجموعہ \(1\) ہونا چاہیے:
\[
K+2K+4K+2K+K=10K=1 \;\Longrightarrow\; K=0.1
\]
تقسیم: \(0.1,\ 0.2,\ 0.4,\ 0.2,\ 0.1\)۔ یہ \(x=4\) کے گرد متناسب ہے، اس لیے
\(E(X)=4\)۔

\textbf{(ج)} \emph{بغیر واپسی} کے نمونہ لینے سے یہ دو ذی حدی نہیں بلکہ
\textenglish{hypergeometric} تقسیم بنتی ہے:
\[
P(X=x)=\frac{\dbinom{4}{x}\dbinom{6}{4-x}}{\dbinom{10}{4}},
\qquad \dbinom{10}{4}=210
\]
\[
P(X=2)=\frac{\dbinom{4}{2}\dbinom{6}{2}}{210}=\frac{6\times 15}{210}
=\frac{90}{210}=\frac{3}{7}\approx 0.4286
\]

\jawab{(الف) \(\tfrac18,\tfrac38,\tfrac38,\tfrac18\)؛ \(P(X\ge2)=0.5\) \quad
(ب) \(K=0.1\) \quad (ج) \(P(X=2)=\tfrac37\approx 0.429\)}
\end{hal}

\begin{ghalti}
حصہ (ج) کو دو ذی حدی مان کر \(p=4/10\) لگانے سے \(0.3456\) آتا ہے --- غلط۔ دو
ذی حدی کے لیے \(p\) کا ہر بار ایک جیسا رہنا ضروری ہے، جو صرف \emph{واپسی کے
ساتھ} ہوتا ہے۔ ''بغیر واپسی'' کے الفاظ دیکھتے ہی
\textenglish{hypergeometric} کی طرف جائیے۔
\end{ghalti}

%---------------------------------------------------------------------
\section*{یونٹ \textenglish{3}: مساوات اور عدم مساوات}
\markboth{یونٹ 3: مساوات اور عدم مساوات}{}

\begin{sawal}{سوال \textenglish{4}}
حل کیجیے: \textbf{(الف)} \(t^{2}+4t=21\) تحلیل سے، \quad
\textbf{(ب)} \(4x^{2}+3x-1=0\) دو درجی فارمولے سے۔
\end{sawal}

\begin{hal}
\textbf{(الف)} پہلے سب کچھ ایک طرف --- تحلیل صرف صفر کے مقابل چلتی ہے۔
\(t^{2}+4t-21=0\)۔ دو عدد جن کا حاصل ضرب \(-21\) اور جمع \(+4\) ہو: \(+7\) اور
\(-3\)۔
\[
(t+7)(t-3)=0 \;\Longrightarrow\; t=-7 \ \text{یا}\ t=3
\]
\textbf{(ب)} \(a=4,\ b=3,\ c=-1\)۔ پہلے ممیز، وہ بتاتا ہے جواب کس قسم کا آئے گا:
\[
b^{2}-4ac=9+16=25>0,\qquad
x=\frac{-3\pm 5}{8} \;\Longrightarrow\; x=\tfrac14 \ \text{یا}\ x=-1
\]
\parakh{\((-7)^{2}+4(-7)=21\)؛ \(4(\tfrac1{16})+3(\tfrac14)-1=0\)۔}

\jawab{(الف) \(t=-7,\ 3\) \quad (ب) \(x=\tfrac14,\ -1\)}
\end{hal}

\begin{sawal}{سوال \textenglish{5}}
حل کیجیے: \textbf{(الف)} \(6t^{2}+t-12>0\)، \quad
\textbf{(ب)} \(|2y+5|=|y-4|\)، \quad \textbf{(ج)} \(|z^{2}-8|\le 8\)۔
\end{sawal}

\begin{hal}
\textbf{(الف)} ممیز \(=1+288=289\)، \(\sqrt{289}=17\):
\[
t=\frac{-1\pm 17}{12} \;\Longrightarrow\; t=\tfrac43 \ \text{یا}\ t=-\tfrac32
\]
\(t^{2}\) کا عامل \(6>0\)، اس لیے قطع مکافی اوپر کھلتا ہے: جڑوں کے درمیان محور
سے نیچے، باہر اوپر۔ ہمیں \(>0\) چاہیے:
\[
t<-\tfrac32 \quad\text{یا}\quad t>\tfrac43
\]
\parakh{\(t=0\) (جڑوں کے بیچ): \(-12\)، یعنی \(>0\) نہیں --- درست طور پر خارج۔}

\textbf{(ب)} \(|u|=|v|\) دو صورتوں میں بٹتا ہے:
\(2y+5=y-4 \Rightarrow y=-9\)؛ اور \(2y+5=-y+4 \Rightarrow y=-\tfrac13\)۔

\textbf{(ج)} \(|u|\le a\) اندر کی طرف بند ہوتا ہے:
\[
-8\le z^{2}-8\le 8 \;\Longrightarrow\; 0\le z^{2}\le 16
\]
بایاں حصہ ہر حقیقی \(z\) کے لیے خود بخود درست ہے؛ دایاں دیتا ہے
\(-4\le z\le 4\)۔

\jawab{(الف) \(t<-\tfrac32\) یا \(t>\tfrac43\) \quad
(ب) \(y=-9,\ -\tfrac13\) \quad (ج) \(-4\le z\le 4\)}
\end{hal}

\begin{ghalti}
\(|u|\le a\) \emph{ایک} وقفہ دیتا ہے؛ \(|u|\ge a\) باہر کھل کر \emph{دو} وقفے
دیتا ہے۔ نیز منفی عدد سے ضرب یا تقسیم پر عدم مساوات کی علامت الٹ جاتی ہے۔
\end{ghalti}

%---------------------------------------------------------------------
\section*{یونٹ \textenglish{4}: خطی مساوات}
\markboth{یونٹ 4: خطی مساوات}{}

\begin{nukta}[پہلے میلان، پھر ایک نقطہ]
\[
m=\frac{y_2-y_1}{x_2-x_1},\qquad y-y_1=m(x-x_1),\qquad y=mx+c
\]
متوازی: \(m_1=m_2\)۔ عمود: \(m_1m_2=-1\)۔ عمودی خط \(x=\)~ثابت کا کوئی میلان
نہیں ہوتا۔
\end{nukta}

\begin{sawal}{سوال \textenglish{6}}
صابن کی طلب \(p=15-\tfrac76 q\) اور رسد \(p=\tfrac23 q\) ہے۔ توازنی مقدار اور
توازنی قیمت نکالیے۔
\end{sawal}

\begin{hal}
منڈی وہاں صاف ہوتی ہے جہاں طلب اور رسد برابر ہوں:
\[
15-\tfrac76 q=\tfrac23 q
\;\xrightarrow{\ \times 6\ }\;
90-7q=4q \;\Longrightarrow\; q^{*}=\tfrac{90}{11}\approx 8.18
\]
\[
p^{*}=\tfrac23\cdot\tfrac{90}{11}=\tfrac{60}{11}\approx 5.45
\]
\parakh{طلب کی طرف سے بھی:
\(15-\tfrac76\cdot\tfrac{90}{11}=15-\tfrac{105}{11}=\tfrac{60}{11}\)۔ دونوں
منحنی متفق ہیں۔}

\jawab{\(q^{*}=\tfrac{90}{11}\approx 8.18\) اکائیاں،
\(p^{*}=\tfrac{60}{11}\approx \$5.45\)}
\end{hal}

\begin{sawal}{سوال \textenglish{7}}
خط کی مساوات نکالیے: \textbf{(الف)} \((-1,3)\) سے گزرے اور \(y=4x-5\) کے متوازی
ہو؛ \textbf{(ب)} \((-5,4)\) سے گزرے اور \(2y=x+1\) پر عمود ہو؛
\textbf{(ج)} \((2,-8)\) سے گزرے اور \(x=-4\) کے متوازی ہو۔
\end{sawal}

\begin{hal}
\textbf{(الف)} متوازی خطوط کا میلان ایک: \(m=4\)۔
\[
y-3=4(x+1) \;\Longrightarrow\; y=4x+7
\]
\textbf{(ب)} \(2y=x+1\) یعنی \(y=\tfrac12 x+\tfrac12\)، میلان \(\tfrac12\)۔
عمود کا میلان منفی معکوس \(=-2\):
\[
y-4=-2(x+5) \;\Longrightarrow\; y=-2x-6
\]
\textbf{(ج)} \(x=-4\) عمودی خط ہے، جس کا میلان ہوتا ہی نہیں۔ اس کے متوازی خط بھی
عمودی ہو گا اور دیے گئے نقطے کی \(x\) قدر سے گزرے گا: \(\ x=2\)۔

\jawab{(الف) \(y=4x+7\) \quad (ب) \(y=-2x-6\) \quad (ج) \(x=2\)}
\end{hal}

\begin{ghalti}
حصہ (ج) اکثر امیدواروں کو پکڑ لیتا ہے۔ \(m_1m_2=-1\) عمودی خطوط پر لاگو نہیں
ہوتا۔ جواب \(x=\)~ثابت ہے، اور ثابت نقطے کی \(x\) قدر سے آتا ہے، \(y\) سے نہیں۔
\end{ghalti}

%---------------------------------------------------------------------
\section*{یونٹ \textenglish{5}: میٹرکس}
\markboth{یونٹ 5: میٹرکس}{}

\begin{sawal}{سوال \textenglish{8}}
\(A=\mat{1&2&3\\-1&-2&3\\3&4&5}\) اور \(B=\mat{4&-1&1\\1&4&3\\3&-3&-2}\) کے لیے
\(AB\) نکالیے، اور تصدیق کیجیے کہ \((A+B)^{t}=A^{t}+B^{t}\)۔
\end{sawal}

\begin{hal}
\(A\) کی سطر کو \(B\) کے ستون سے ضرب دیجیے۔ پہلی سطر \((1,2,3)\) کے لیے:
\[
4+2+9=15,\qquad -1+8-9=-2,\qquad 1+6-6=1
\]
اسی طرح باقی سطریں:
\[
AB=\mat{15&-2&1\\3&-16&-13\\31&-2&5}
\]
\textbf{تصدیق:} \(A+B=\mat{5&1&4\\0&2&6\\6&1&3}\) لہٰذا
\((A+B)^{t}=\mat{5&0&6\\1&2&1\\4&6&3}\)، اور
\[
A^{t}+B^{t}=\mat{1&-1&3\\2&-2&4\\3&3&5}+\mat{4&1&3\\-1&4&-3\\1&3&-2}
=\mat{5&0&6\\1&2&1\\4&6&3}
\]
دونوں برابر ہیں۔

\jawab{\(AB=\mat{15&-2&1\\3&-16&-13\\31&-2&5}\)؛ تصدیق مکمل۔}
\end{hal}

\begin{ghalti}
\((AB)^{t}=B^{t}A^{t}\) ہے، \(A^{t}B^{t}\) نہیں --- اور اسی طرح
\((AB)^{-1}=B^{-1}A^{-1}\)۔ دونوں میں ترتیب الٹ جاتی ہے۔ یہ رواج نہیں، ابعاد
کے میل کی مجبوری ہے۔
\end{ghalti}

%---------------------------------------------------------------------
\section*{یونٹ \textenglish{6}: محدد اور معکوس}
\markboth{یونٹ 6: محدد اور معکوس}{}

\begin{nukta}[ہم عامل، معکوس، کریمر]
\(C_{ij}=(-1)^{i+j}M_{ij}\)؛ علامتوں کا نقشہ \(\mat{+&-&+\\-&+&-\\+&-&+}\)۔
\[
\det A=\sum_j a_{ij}C_{ij},\qquad
\operatorname{adj}A=\bigl[C_{ij}\bigr]^{t},\qquad
A^{-1}=\frac{\operatorname{adj}A}{\det A}
\]
\textbf{کریمر:} \(x_k=\dfrac{D_k}{D}\)، جہاں \(D_k\) میں \(k\)واں ستون \(b\) سے
بدلا جاتا ہے۔ شرط \(D\neq0\)۔ سب سے زیادہ صفر والی سطر سے پھیلائیے؛ دو یکساں
سطریں ہوں تو \(\det=0\)۔
\end{nukta}

\begin{sawal}{سوال \textenglish{9}}
محدد نکالیے: \quad \textbf{(الف)} \(\dt{2&3&-1\\1&1&0\\2&-3&5}\) \qquad
\textbf{(ب)} \(\dt{2a&a&a\\b&2b&b\\c&c&2c}\)
\end{sawal}

\begin{hal}
\textbf{(الف)} دوسری سطر میں ایک صفر ہے، اسی سے پھیلائیے:
\[
2\dt{1&0\\-3&5}-3\dt{1&0\\2&5}-\dt{1&1\\2&-3}=2(5)-3(5)-(-5)=0
\]
محدد کا صفر ہونا بتاتا ہے کہ سطریں آپس میں وابستہ ہیں۔

\textbf{(ب)} ہر سطر سے مشترک عامل باہر نکالیے:
\[
abc\dt{2&1&1\\1&2&1\\1&1&2}=abc\bigl[2(3)-1(1)+1(-1)\bigr]=4abc
\]

\jawab{(الف) \(0\) \quad (ب) \(4abc\)}
\end{hal}

\begin{sawal}{سوال \textenglish{10}}
\(A=\mat{1&2&3\\1&3&5\\1&5&12}\) کا معکوس ہم عاملوں کے طریقے سے نکالیے۔
\end{sawal}

\begin{hal}
\[
\det A=1(36-25)-2(12-5)+3(5-3)=11-14+6=3\neq0
\]
ہم عامل نکال کر \emph{منقلب} لیجیے --- یہی ملحقہ ہے:
\[
\operatorname{adj}A=\mat{11&-9&1\\-7&9&-2\\2&-3&1},
\qquad
A^{-1}=\frac{1}{3}\mat{11&-9&1\\-7&9&-2\\2&-3&1}
\]
\parakh{\(A\) کی پہلی سطر \((1,2,3)\) ضرب \(A^{-1}\) کا پہلا ستون
\(\tfrac13(11,-7,2)^{t}\): \(\tfrac{11-14+6}{3}=1\)۔ تیسری سطر \((1,5,12)\) اسی
ستون سے: \(\tfrac{11-35+24}{3}=0\)۔}

\jawab{\(A^{-1}=\dfrac{1}{3}\mat{11&-9&1\\-7&9&-2\\2&-3&1}\)}
\end{hal}

\begin{ghalti}
ملحقہ ہم عاملوں کے میٹرکس کا \textbf{منقلب} ہے۔ منقلب بھول جانے سے
\(AA^{-1}=I\) خاص طور پر \emph{غیر قطری} خانوں میں ٹوٹتا ہے --- اسی لیے پڑتال
ایک غیر قطری خانے پر بھی کیجیے، صرف قطری پر نہیں۔
\end{ghalti}

\begin{sawal}{سوال \textenglish{11}}
کریمر کے اصول سے حل کیجیے: \quad
\(x+y=2\)، \quad \(2x-z=1\)، \quad \(2y-3z=-1\)
\end{sawal}

\begin{hal}
سب سے پہلے غائب متغیرات کے صفر عامل \emph{لکھ} لیجیے، ورنہ اعداد غلط ستون میں
چلے جائیں گے:
\[
x+y+0z=2,\qquad 2x+0y-z=1,\qquad 0x+2y-3z=-1
\]
\[
D=\dt{1&1&0\\2&0&-1\\0&2&-3}=1(0+2)-1(-6-0)+0=8
\]
\[
D_x=\dt{2&1&0\\1&0&-1\\-1&2&-3}=8,\quad
D_y=\dt{1&2&0\\2&1&-1\\0&-1&-3}=8,\quad
D_z=\dt{1&1&2\\2&0&1\\0&2&-1}=8
\]
\[
x=\tfrac88=1,\qquad y=1,\qquad z=1
\]
\parakh{\(1+1=2\)؛ \(2(1)-1=1\)؛ \(2(1)-3(1)=-1\)۔ تینوں درست۔}

\jawab{\(x=1\)، \(y=1\)، \(z=1\)}
\end{hal}

\begin{ghalti}
کریمر لگانے سے \emph{پہلے} ہمیشہ \(D\) نکالیے۔ \(D=0\) ہو تو یہ اصول لاگو نہیں
ہوتا، اور درست جواب اعداد نہیں بلکہ نظام کے بارے میں بیان ہوتا ہے: ''کوئی حل
نہیں'' یا ''لامتناہی حل''۔
\end{ghalti}

%---------------------------------------------------------------------
\section*{یونٹ \textenglish{7}: حد اور تفرق}
\markboth{یونٹ 7: حد اور تفرق}{}

\begin{sawal}{سوال \textenglish{12}}
حد نکالیے: \quad
\textbf{(الف)} \(\displaystyle\lim_{x\to 3}\frac{x^{4}-81}{x^{2}-x-6}\) \qquad
\textbf{(ب)} \(\displaystyle\lim_{x\to \infty}\frac{4-3x^{3}}{x^{3}-1}\)
\end{sawal}

\begin{hal}
\textbf{(الف)} تعویض سے \(\tfrac00\) آتا ہے، لہٰذا تحلیل کیجیے:
\[
x^{4}-81=(x-3)(x+3)(x^{2}+9),\qquad x^{2}-x-6=(x-3)(x+2)
\]
مشترک عامل \((x-3)\) منسوخ:
\[
\lim_{x\to3}\frac{(x+3)(x^{2}+9)}{x+2}=\frac{(6)(18)}{5}=\frac{108}{5}=21.6
\]
\textbf{(ب)} اوپر اور نیچے دونوں کو \(x^{3}\) سے تقسیم کیجیے:
\[
\lim_{x\to\infty}\frac{\tfrac{4}{x^{3}}-3}{1-\tfrac{1}{x^{3}}}=\frac{-3}{1}=-3
\]
مختصر راستہ: برابر درجے کے کثیر رقمی کسر کی حد سرکردہ عاملوں کا تناسب ہوتی ہے۔

\jawab{(الف) \(\tfrac{108}{5}=21.6\) \quad (ب) \(-3\)}
\end{hal}

\begin{sawal}{سوال \textenglish{13}}
\(x\) اکائیوں کی لاگت \(C(x)=x^{2}-5x+80\) ہے۔ \textbf{(الف)} \textenglish{10}
اور \textenglish{15} کے درمیان تبدیلی کی اوسط شرح، اور \textbf{(ب)}
\textenglish{12} اکائیوں پر آنی شرح نکالیے۔
\end{sawal}

\begin{hal}
\textbf{(الف)} اوسط شرح دو نقطوں کو ملانے والے وتر کا میلان ہے:
\[
C(15)=230,\quad C(10)=130
\;\Longrightarrow\;
\frac{230-130}{15-10}=\frac{100}{5}=20
\]
\textbf{(ب)} آنی شرح مشتق ہے --- یہی \emph{حاشیائی لاگت} ہے:
\[
C'(x)=2x-5 \;\Longrightarrow\; C'(12)=24-5=19
\]

\jawab{(الف) \(20\) فی اکائی \quad (ب) \(19\) فی اکائی}
\end{hal}

\begin{ghalti}
''اوسط شرح'' کے لیے دو نقطے اور تفریق چاہیے؛ ''آنی شرح'' کے لیے ایک نقطہ اور
مشتق۔ یہ الگ مقداریں ہیں اور یہاں الگ جواب دیتی ہیں (\(20\) بمقابلہ \(19\))۔
\end{ghalti}

%---------------------------------------------------------------------
\section*{یونٹ \textenglish{8}: جزوی مشتق}
\markboth{یونٹ 8: جزوی مشتق}{}

\begin{sawal}{سوال \textenglish{14}}
ایک کمپنی کا منافع \(P(x,y)=1200+80x-2x^{2}+100y-y^{2}\) ہے، جہاں \(x\) محنت اور
\(y\) اشیا کی فی اکائی لاگت ہے۔ منافع کو زیادہ سے زیادہ کرنے والے \(x\)، \(y\)
اور زیادہ سے زیادہ منافع نکالیے۔
\end{sawal}

\begin{hal}
\(xy\) کی رقم موجود نہیں، اس لیے دونوں متغیر الگ الگ بہتر ہوتے ہیں:
\[
\pd{P}{x}=80-4x=0 \Rightarrow x=20,
\qquad
\pd{P}{y}=100-2y=0 \Rightarrow y=50
\]
\textbf{دوسرے مشتق کی جانچ:} \(P_{xx}=-4\)، \(P_{yy}=-2\)، \(P_{xy}=0\)، لہٰذا
\[
D=(-4)(-2)-0=8>0 \ \text{اور}\ P_{xx}<0 \;\Longrightarrow\; \text{اعظم}
\]
\[
P(20,50)=1200+1600-800+5000-2500=4500
\]
\parakh{کسی بھی متغیر کو ہلانے سے منافع گھٹتا ہے: \(P(21,50)=4498\)،
\(P(20,51)=4499\)۔}

\jawab{\(x=20\)، \(y=50\)؛ زیادہ سے زیادہ منافع \textenglish{4{,}500}۔}
\end{hal}

%---------------------------------------------------------------------
\section*{یونٹ \textenglish{9}: بہتر بندی}
\markboth{یونٹ 9: بہتر بندی}{}

\begin{nukta}[کاروبار کے چار فعل]
\[
R(x)=x\,p(x),\qquad P(x)=R(x)-C(x),\qquad \bar{C}(x)=\frac{C(x)}{x}
\]
حاشیائی روپ مشتق ہیں: \(MR=R'\)، \(MC=C'\)۔ منافع وہاں اعظم ہوتا ہے جہاں
\(P'=0\) یعنی \(MR=MC\)، بشرطیکہ \(P''<0\) ہو۔
\end{nukta}

\begin{sawal}{سوال \textenglish{15}}
کل لاگت \(c=0.03x^{2}+8x+300\) ہے۔ کس پیداوار پر فی اکائی اوسط لاگت کم سے کم ہو
گی؟
\end{sawal}

\begin{hal}
\[
\bar{c}(x)=\frac{0.03x^{2}+8x+300}{x}=0.03x+8+\frac{300}{x}
\]
\[
\bar{c}\,'(x)=0.03-\frac{300}{x^{2}}=0
\;\Longrightarrow\;
x^{2}=10{,}000 \;\Longrightarrow\; x=100
\]
(مثبت جڑ، کیونکہ پیداوار منفی نہیں ہو سکتی۔)
\(\bar{c}\,''=\tfrac{600}{x^{3}}>0\)، لہٰذا یہ اصغر ہے۔
\[
\bar{c}(100)=3+8+3=14
\]
\parakh{اصغر پر حاشیائی لاگت اوسط لاگت کے برابر ہوتی ہے:
\(c'(100)=0.06(100)+8=14\)۔ یہ ہمیشہ سچ ہے اور بہترین پڑتال ہے۔}

\jawab{\textenglish{100} اکائیوں پر، جہاں اوسط لاگت \textenglish{\$14} فی اکائی
ہے۔}
\end{hal}

\begin{sawal}{سوال \textenglish{16}}
طلب \(p=52-0.02x\) اور لاگت \(c(x)=400+40x\) ہے۔ کس پیداوار پر منافع اعظم ہو گا؟
قیمت اور منافع کتنا ہو گا؟
\end{sawal}

\begin{hal}
\[
R(x)=x(52-0.02x)=52x-0.02x^{2}
\]
\[
P(x)=R-c=52x-0.02x^{2}-400-40x=12x-0.02x^{2}-400
\]
\[
P'(x)=12-0.04x=0 \Rightarrow x=300,
\qquad P''=-0.04<0 \ (\text{اعظم})
\]
\[
p=52-0.02(300)=46,
\qquad
P(300)=3600-1800-400=1400
\]
\parakh{\(MR=MC\) سے بھی: \(52-0.04x=40\) دیتا ہے \(x=300\) --- وہی پیداوار، مگر
معاشی اصول سے۔}

\jawab{پیداوار \textenglish{300} اکائیاں، قیمت \textenglish{\$46}، منافع
\textenglish{\$1{,}400}۔}
\end{hal}

\begin{ghalti}
اعظم منافع اور اعظم آمدنی ایک چیز نہیں۔ یہاں آمدنی \(x=1300\) پر عروج پر ہوتی
ہے، جو منافع بخش \textenglish{300} سے بہت آگے ہے۔ اسی فعل کو بہتر کیجیے جس کا
نام سوال میں ہو۔
\end{ghalti}
